\par\medskip
\SourceLeaf{078}\pdfbookmark[2]{59.}{article-059}
\begin{center}59.\end{center}
Praecepta pro determinandis distantiis a planis per Solem transeuntibus manifesto etiam ad distantias terrae applicare licet: hic vero simplicissimi tantum casus occurrere solent. Sit $R$ distantia terrae a Sole, $L$ longitudo heliocentrica terrae (quae $180^0$ a longitudine geocentrica Solis differt), denique $X$, $Y$, $Z$ distantiae terrae a tribus planis in Sole sub angulis rectis se secantibus. Iam si

I. Planum ipsarum $Z$ est ipsa ecliptica, longitudinesque polorum planorum reliquorum, a quibus distantiae sunt $X$, $Y$, resp. $N$ et $N+90^0$: erit
\[
X=R\cos(L-N),\qquad Y=R\sin(L-N),\qquad Z=0.
\]

II. Si planum ipsarum $Z$ aequatori parallelum est, atque rectascensiones polorum planorum reliquorum, a quibus distantiae sunt $X$, $Y$, resp. $0$ et $90^0$, habebimus, obliquitate eclipticae per $\epsilon$ designata,
\[
X=R\cos L,\qquad Y=R\cos\epsilon\sin L,\qquad Z=R\sin\epsilon\sin L.
\]
Tabularum solarium recentissimarum editores, clarr. \textsc{de Zach} et \textsc{de Lambre}, latitudinis Solis rationem habere coeperunt, quae quantitas a perturbationibus reliquorum planetarum atque lunae producta vix unum minutum secundum attingere potest. Designando latitudinem heliocentricam terrae, quae latitudini Solis semper aequalis sed signo opposito affecta erit, per $B$, habebimus:
\[
\begin{array}{@{}l|l@{}}
\text{in casu I.}&\text{in casu II.}\\[.4em]
X=R\cos B\cos(L-N)&X=R\cos B\cos L\\[.4em]
Y=R\cos B\sin(L-N)&Y=R\cos B\cos\epsilon\sin L-R\sin B\sin\epsilon\\[.4em]
Z=R\sin B&Z=R\cos B\sin\epsilon\sin L+R\sin B\cos\epsilon.
\end{array}
\]
Pro $\cos B$ hic semper tuto substitui poterit $1$, angulusque $B$ in partibus radii expressus pro $\sin B$.

Coordinatae ita inventae ad \textit{centrum} terrae referuntur: si $\xi$, $\eta$, $\zeta$ sunt distantiae puncti cuiuslibet in terrae \textit{superficie} a tribus planis per centrum terrae ductis iisque quae per Solem ducta erant parallelis, distantiae illius puncti a \SourceLeaf{079}planis per Solem transeuntibus manifesto erunt $X+\xi$, $Y+\eta$, $Z+\zeta$, valores coordinatarum $\xi$, $\eta$, $\zeta$ autem pro utroque casu facile determinantur sequenti modo. Sit $\rho$ radius globi terrestris (sive sinus parallaxis horizontalis mediae Solis), $\lambda$ longitudo puncti sphaerae coelestis, ubi recta a terrae centro ad punctum superficiei ductum proiicitur, $\beta$ eiusdem latitudo, $\alpha$ ascensio recta, $\delta$ declinatio, eritque
\[
\begin{array}{@{}l|l@{}}
\text{in casu I.}&\text{in casu II.}\\[.4em]
\xi=\rho\cos\beta\cos(\lambda-N)&\xi=\rho\cos\delta\cos\alpha\\[.4em]
\eta=\rho\cos\beta\sin(\lambda-N)&\eta=\rho\cos\delta\sin\alpha\\[.4em]
\zeta=\rho\sin\beta&\zeta=\rho\sin\delta.
\end{array}
\]
Punctum illud sphaerae coelestis manifesto respondet ipsi zenith loci in superficie (siquidem terra tamquam sphaera spectatur), quocirca ipsius ascensio recta conveniet cum ascensione recta medii coeli sive cum tempore siderali in gradus converso, declinatio autem cum elevatione poli; si operae pretium esset, figurae terrestris sphaeroidicae rationem habere, pro $\delta$ elevationem poli \textit{correctam}, atque pro $\rho$ distantiam veram loci a centro terrae accipere oporteret, quae per regulas notas eruuntur. Ex $\alpha$ et $\delta$ longitudo et latitudo $\lambda$ et $\beta$ per regulas notas infra quoque tradendas deducentur; ceterum patet, $\lambda$ convenire cum longitudine \textit{nonagesimi}, atque $90^0-\beta$ cum eiusdem altitudine.
